\section{引言}
\label{sec:introduction}

\subsection{城市犯罪计数预测中的数据价值问题}
\label{subsec:intro_data_value}

城市犯罪预测越来越多地把历史事件与移动轨迹、日历变量、环境测量和其他动态城市
数据结合。这些数据能够丰富对地点周边活动的描述，但也需要持续投入采集、清洗、
空间对齐、维护和治理成本。因此，对城市管理部门或分析人员而言，关键问题不仅是
加入新数据后某个拟合模型的测试误差是否下降，还包括该数据是否在地点、犯罪历史
和日历信息之外进一步降低了未来犯罪计数的不确定性。

这两个问题并不等价。样本外误差的变化取决于拟合模型、超参数搜索、训练样本、目标
定义和验证时期，既可能来自有效的新信息，也可能来自新增变量与特定模型类别的匹配
或抽样波动。相比之下，模型无关预测基准由目标与发布预测时可用信息的联合分布决定。
把实际误差与该基准比较，有助于区分目标--信息对的内在不确定性和仍可能通过改进
建模而降低的误差。

犯罪计数预测使这种区分具有直接的应用意义。事件计数是非负整数，在精细时空尺度上
通常含有大量零值，并常以均方误差（MSE）等损失进行评价。现有研究已发展统计、
机器学习、循环、图结构、多尺度和 Transformer 等模型
\citep{kounadi2020review,du2023review,mandalapu2023review,
alves2018crime,angbera2023model,huang2018deepcrime,deng2023risk,
jing2024multiscale,wang2025mragnn,butt2025start,cui2026transfer,
hu2023fusion}。这些比较对模型选择非常重要，却不能识别在指定信息集下有多少误差
无法避免。当一个成本较高的新城市数据源看似改善一个或多个模型时，这一局限尤其
重要。

本文以共享自行车活动作为候选新增信息源研究上述问题。自行车行程提供一种高频、
空间细化的公共空间活动观测，但并不构成 ambient population 普查；其与犯罪之间
的预测关联也不等于因果效应。本文提出的问题更为严格和有限：滞后自行车流是否为
未来犯罪计数提供额外预测信息，以及估计出的增益能否区别于有限样本伪象？

\subsection{从整数值不确定性到有限样本校准证据}
\label{subsec:intro_calibrated_evidence}

条件熵自然衡量观测预测时信息集后仍然保留的不确定性
\citep{shannon1948,cover_thomas}。然而，现有基于熵的预测下界主要面向连续结果
或分类损失。连续高斯熵--方差论证不能通过简单地把微分熵替换为 Shannon 熵而应用
于整数犯罪计数。预测误差位于整数格点上，二阶矩约束下对应的最大熵包络是离散的。

为解决这一支撑不匹配问题，本文建立离散熵下界（DELB）。DELB 通过反演整数格点
最大熵包络，把条件 Shannon 熵映射为因果整数值预测器的模型无关、熵约束 MSE
底线。该下界在离散包络形式上是精确的，但只有满足给定可达条件时才等于最优整数
预测器风险。框架还给出预测器相关强化，并区分两种不可达来源：预测误差对可用历史
仍存在依赖，以及误差分布与匹配二阶矩的离散高斯之间存在形状失配。

总体排序只是应用问题的第一层。加入离散化移动变量会进一步划分基准状态空间，通常
产生稀疏或 singleton 单元。条件熵与条件互信息（CMI）估计量在基准状态表和增强
状态表中可能具有不同的有限样本误差
\citep{paninski2003,calcagnile2026benchmark}。因此，即使新增变量在总体中没有
信息价值，也可能出现正 CMI 估计或较低 DELB 估计。此外，随机化参考依赖其具体
变换设计；当设计改变支持占用时，其中心不必为零
\citep{ernst2004permutation,zan2022cmi,runge2018cmi,li2023nncit}。

评价过程因此必须区分四个对象：总体 DELB 与 CMI estimand、有限样本估计、相对于
预设随机化参考的证据，以及实际样本外预测误差。本文利用已知真值模拟、支持度诊断、
区组重采样、时间与空间随机化、循环移位参考和匹配预测模型连接这些层次。这一设计
不会把变换特异诊断变成因果检验或精确条件独立检验；其目的是判断观测统计量能否
区别于各项明确参考设计下的有限样本行为。

\subsection{研究问题与贡献}
\label{subsec:intro_contributions}

本文回答三个相互关联的问题：

\begin{enumerate}
    \item \textit{如何把条件不确定性转换为因果整数值预测的离散形式、模型无关
    MSE 下界？}
    \item \textit{如何把估计信息增益与稀疏状态估计偏差及变换特异随机化参考
    区分？}
    \item \textit{有限样本校准后，滞后自行车移动是否仍为城市犯罪计数预测提供
    可区分的额外信息？}
\end{enumerate}

对应贡献强调应用价值。第一，本文把 DELB 建立为整数值城市计数预测的模型无关
基准，给出可达性诊断及可实施的数值反演程序。第二，本文提出有限样本校准流程，
区分估计信息增益、支持结构诱发的偏差、设计条件随机化证据和实际预测改善。第三，
本文利用 2020--2022 年华盛顿特区、纽约和温哥华的统一犯罪与共享自行车数据开展
多城市评价，考察城市与局部估计、时空尺度、样本量、替代估计器及匹配样本外预测。

经验结论保持审慎。加入滞后自行车流降低了主要城市比较中的 DELB 点估计，但校准后
这些降幅无法区别于预设的时间、空间或循环移位参考，实际预测变化也依赖模型和城市。
本文不把这一结果视为应用失败，而是用它阐明：评价新型智能城市数据时，较低估计
不确定性、经校准的信息增益和预测性能改善必须作为不同主张分别报告。

第~\ref{sec:related_work} 节综述城市犯罪预测、犯罪分析中的移动数据及信息论预测
下界；第~\ref{sec:theory} 节给出 DELB 方法；第~\ref{sec:data} 节定义多城市数据
与实验设计；第~\ref{sec:results} 节先汇总多城市应用结论，再报告有限样本、城市、
多尺度和预测证据；
第~\ref{sec:discussion} 节讨论实际意义、推断范围与局限。
